\section{课题来源及研究的目的和意义}\label{Sec: purpose and background}
\vspace{0.8em}

\vspace{0.8em}
\subsection{研究背景}\label{background}
\vspace{0.8em}



深圳出台《深圳市低空经济产业创新发展实施方案（2022-2025年）》《深圳市支持低空经济高质量发展的若干措施》《深圳经济特区低空经济产业促进条例》等，从政策、法律等多方面为深圳先行先试制定民用无人机管理规则和运行标准、推动低空经济发展提供有力支撑。低空经济的发展重点是使用小型无人机在城市低空环境进行社会服务工作。在民生领域，美团、顺丰等高科技公司，开辟低空运输的新型配送方式；在农业领域，小型无人机被广泛用于喷洒农药、播撒种子等作业，极大地提高了农事活动的自动化与精细化水平；此外，配备摄像设备的小型无人机还广泛应用于高空航拍、环境监测以及新闻报道等领域，充分展示了其在推动经济发展，培育发展新质生产力的重要作用。

在城市密集区域，建筑物及基础设施引发的湍流和阵风会导致局部空气流场剧烈变化，进而干扰无人机周围流场，造成飞行抖动、轨迹偏移，从而增加与人员、建筑物甚至其他无人机发生碰撞的风险，严重威胁低空运行的安全性。








\vspace{0.8em}
\subsection{研究的目的及意义}\label{purpose}
\vspace{0.8em}

螺旋桨是为无人机提供升力、保证无人机能够正常起飞的关键组成部分。由于小型无人机在尺寸上比普通直升飞机等大大缩小，其螺旋桨的空气动力学性能规律与普通直升飞机的不完全相同。因而需要针对小型无人机的螺旋桨研究其空气动力学性能，进而对其螺旋桨进行气动外形优化设计，以减小飞行过程中的阻力，提高飞行过程中的升力，进而起到提升小型无人机的飞行效率，实现在相同的电池供能的条件下，无人机能够增大飞行距离或运输更重的物品。%根据老板目标再做修改。

螺旋桨的气动外形优化问题需要获取螺旋桨翼型的空气动力学参数，应用合适的优化算法，寻到最优解。在整个优化过程中，空气动力学参数的获取难度大，耗时长，空气动力学性能参数的获取方法主要有风洞实验法，有高保真度的计算流体力学（Computational fluid dynamics, CFD）方法，以及随着近年人工智能领域发展而发展的代理模型法。
风洞实验是研究螺旋桨空气动力学性能的方法之一，它能够获得高精度的数据，可信度高。但是风洞实验所需的人力物力成本高，因此国内外的研究人员已不再将其列为研究方法的首选。取而代之的是具有高保真度的CFD方法。在使用CFD方法的气动外形优化的过程中，常常需要对螺旋桨的几何尺寸进行调整变动，每次几何尺寸的改变，都需要对新尺寸的螺旋桨的流场进行计算，如果采用复杂的三维流场作为研究对象，则计算耗时长，难以适应不断变化的几何调整方案。为了能够避免复杂的CFD计算，减少计算时间，保证优化任务较快完成，研究人员主要采取一种简化的方案。
该简化方案采用精度不高但计算速度快的方法，快速获取气动特性参数。当得到优化结果后再进行高精度的CFD计算，验证优化结果的可信度。随着人工智能的发展，近年来研究人员多采用代理模型来代替耗时的CFD计算，从而实现高效的气动外形设计研究。在构建代理模型之前，需要进行大量CFD计算或收集实验数据，以确保代理模型的准确，常用的代理模型有神经网络、本征正交分解等。
本课题旨在通过建立多重保真度模型，搭建低保真度数据与高保真度数据的桥梁，从而实现通过快速获取的低保真数据和多保真度模型，即可得到难获取的高保真度的数据，提高数据的精度，以缩短整个螺旋桨气动外形优化问题的时间。

优化问题的关键是选择合适的优化方法，好的优化方法能够主导整个优化过程的发展方向。一般来说，优化方法可以分为基于梯度的优化方法和无梯度优化方法。基于梯度的优化方法需要求解梯度，因而需要较长时间，且容易陷入局部最优解，难以获得全局最优解。基于梯度的优化方法主要有最速下降法（The steepest descent method）、内点法（Interior point method）等。无梯度优化方法因其不需要计算梯度，降低了计算的复杂度和开发难度。启发式算法是应用最广泛的一类无梯度优化方法，如遗传算法（Genetic algorithm, GA）、粒子群算法（Particle swarm optimization, PSO）、模拟退火算法（Simulated annealing, SA）等。优化方法的选择较难，往往凭借研究人员的经验或采用不同的方法进行对比，因此需要一种能够指导选择优化方法的手段，以帮助选择合适的优化方法，进而寻到计算域内的最优解。本课题旨在探索不同优化方法的适用条件，从而指导优化问题的优化方法的选择。

本课题拟采用数据拓扑结构帮助选择合适的优化方法。数据的拓扑特征是帮助分析数据的有利方法，数据的拓扑特征能够描述数据及其响应面的形状，了解数据点与数据点之间的关系，进而帮助分析数据的特点，挖掘隐藏在数据结构中的信息，从而选择合适的优化方法解决优化问题。拓扑特征是指描述数据集拓扑结构的属性，它包括极值的个数、边界点、数据山脊线等。比如，数据集中只包含一个极值点，那么梯度下降法即可找到最优解。如果在高维数据域中有几个极值点，则需要使用随机梯度下降法。如果十多个极值点时，则需要使用如遗传算法的方法。

总之，本课题的研究目的是提高小型四旋翼无人机螺旋桨的空气动力学性能，基于数据的拓扑结构分析，给出一种指导选择优化方法的手段，采用一种具有多重保真度模型，将低保真度的数据精度提高，缩短获得高精度数据的时间，从而加速螺旋桨气动外形优化设计的过程。